% Source: sources/returns/S6_04_MAPS_MANIFOLDS/code/manifolds_invariant_chern_caps_proof.tex:1–262
\paragraph{Theorem D.2.51 (Proved: invariant-Chern original additive caps and their integral extension).}
Let $K$ be the reduced central fibre of an original additive
Kodaira neighbourhood, of type
$II,III,IV,IV^*,III^*,I_1^*$ or $I_3^*$.
Write its rational components as $C_i$, its primitive
multiplicity vector as $\mu=(\mu_i)$, and let $\sigma$
be the actual component permutation induced by the section $Q$.
Let the original line bundle have invariant component degrees
$c_i$, so that
\[
                  \sigma c=c,\qquad \sum_i\mu_i c_i=0.
                                                        \tag{D.IC1}
\]
These include every corrected line bundle in D.2.50.
Let the component orbits be $O_1,\ldots,O_\rho$, with
lengths $l_j$ and common degrees $c_j$ on $O_j$.
For $c\ne0$ put
\[
                  g=\gcd_j(c_j),\qquad
                  h=\gcd_j(l_jc_j),
                                                        \tag{D.IC2}
\]
where both gcds are positive. For the resulting original
six-dimensional cap $U$, there is an exact integral cellular
cone, with its actual circle-fibre and regular-area classes,
\[
\begin{gathered}
 C_0=\mathbb Z,\quad C_1=\mathbb Z^2,\quad
 C_2=\mathbb Z^v\oplus\mathbb Zt,\quad
 C_3=\mathbb Z^v\oplus\mathbb Z^v,\quad C_4=\mathbb Z^v,\\
 d_2=\begin{pmatrix}c^t&0\\0&0\end{pmatrix},\qquad
 d_3=\begin{pmatrix}0&\sigma-I\\0&c^t\end{pmatrix},\qquad
 d_4=\begin{pmatrix}I-\sigma\\0\end{pmatrix}.
\end{gathered}                                             \tag{D.IC3}
\]
The degree-one generators are the scalar circle $\delta$
and a positive long circle $\gamma$. In particular
\[
\begin{array}{c|ccccc}
i&0&1&2&3&4\\\hline
H_i(U;\mathbb Z)&\mathbb Z&
\mathbb Z\oplus\mathbb Z/g&
\mathbb Z^{\rho-1}\oplus\mathbb Z/h&
\mathbb Z^{2\rho-1}&\mathbb Z^\rho .
\end{array}                                               \tag{D.IC4}
\]
Moreover $\pi_1(U)=\mathbb Z\oplus\mathbb Z/g$.
When $c=0$, these groups are instead
\[
 (H_0,H_1,H_2,H_3,H_4)=
 (\mathbb Z,\mathbb Z^2,\mathbb Z^{\rho+1},
                         \mathbb Z^{2\rho},\mathbb Z^\rho).
                                                        \tag{D.IC5}
\]
The actual regular-fibre torus $\delta\wedge\gamma$
maps to the class of $t$ in (D.IC3), which generates
\emph{all} of the cyclic torsion $\mathbb Z/h$ when
$c\ne0$. The elliptic tori crossed with $\delta$ map
to zero. The complex elliptic area maps to the literal
cycle $(\mu,0)$ in (D.IC3).

Let $B$ be the elliptic monodromy, $q$ its marked
$Q$-cocycle, and $z=(B-I)^{-1}q$. Put
\[
        q_Q=\gcd_j(l_j\mu_j)
             =\operatorname{ord}(\overline Q).
                                                        \tag{D.IC6}
\]
The actual logarithm $z$ has exact denominator $q_Q$.
For $c\ne0$, the full regular-fibre image in $H_2(U)$
is $\mathbb Z\oplus\mathbb Z/h$ with free cokernel
of rank $\rho-2$. For $c=0$ it is a primitive
$\mathbb Z^2$ with free cokernel of rank $\rho-1$.
In both cases this is already the full boundary image.
Consequently the full boundary kernel projects onto
all of $\operatorname{Fix}(A)$ in the Wang sequence;
there is no omitted finite normal index.

If $h=1$, these facts determine the complete marked
local kernels. In the regular basis
$(e_1,e_2,\delta,\gamma)$ the actual fibre kernel is
\[
 \ker\bigl(q_Q(1,0,z_2,0,-z_1,0):\mathbb Z^6\to\mathbb Z\bigr),
                                                        \tag{D.IC7}
\]
where the six coordinates have their usual increasing
exterior order. Both literal fixed-normal tori bound.
The cap kills every vertical fundamental-group
coordinate and retains the primitive $\gamma$.
This includes the canonical new corrected star degrees
\[
 \begin{array}{c|c}
 I_1^*&(1,1,0,0,-1,0)\\
 I_3^*&(1,1,0,0,-1,0,0,0)
 \end{array}
                                                        \tag{D.IC8}
\]
for component classes $P=Q=2$ in $\mathbb Z/4$.
Their actual permutation has two paired outer orbits
and fixed inner components, and $g=h=1$.

For $h>1$, (D.IC3)--(D.IC6) do not by themselves specify
the torsion corrections of the two mixed
elliptic--gamma tori and the second fixed-normal torus.
The theorem does not replace those marked columns by
their saturation. In particular it is not yet an
unrestricted global gluing algorithm for those caps.

\begin{proof}
\emph{The actual original cap and the circle bundle.}
The reduced additive fibre is simply connected and
is homotopy equivalent to a wedge of its $v$ rational
component spheres. The tangency in type $III$ and the
triple point in type $IV$ do not add a graph loop;
the same assertion follows directly by taking the
union of the rational components and their actual
connected intersections. The other listed fibres
have a tree of rational components.
The elliptic neighbourhood retracts onto this fibre.
The unit original linearization makes its compact
cap the mapping torus of the principal scalar-circle
bundle $P_c\to K$, with the lift of $Q$.
This is the same actual original filling used in
D.2.31--32, now retaining the full degree vector.

A principal circle bundle over this wedge is determined
by its $v$ Chern evaluations. Its cellular complex has
one zero-cell, the scalar one-cell, $v$ two-cells with
boundary degrees $c_i$, and $v$ three-cells with zero
boundary. The lift of $Q$ acts by the component
permutation on the two- and three-cells. There is no
additional homotopy class of a gauge transformation:
$[K,S^1]=H^1(K;\mathbb Z)=0$. More explicitly, after
choosing a fixed point of the induced component-tree
map and a fibre point there, the component maps are
homotopic to their orientation-preserving permutations.
The invariant vector $c$ lifts this homotopy to the
circle bundle. A comparison unit has a continuous
logarithm on $K$, so its lift is homotopic through
bundle maps to the identity. A constant scalar at
the chosen point is removed by a scalar isotopy.
This fixes the scalar circle and its product with
the mapping-torus direction throughout the comparison.
The signs in (D.IC3) use the literal product orientation
$t=\delta\wedge\gamma$: the boundary of a two-cell
crossed with the long interval contributes $+c_i t$.
This is the usual mapping cone after changing the signs
of its shifted summands. Thus (D.IC3) is an actual model, not
merely a rational model with the same permutation.

The circle-bundle homotopy sequence gives
$\pi_1(P_c)=\mathbb Z/g$ for $c\ne0$.
The return fixes the scalar generator; the mapping-torus
extension by the infinite cyclic group therefore splits.
This proves the fundamental-group statement and $H_1$.

\emph{The integral extension.}
In $C_2/\operatorname{im}d_3$, successive cells in an
orbit have the relations
\[
                  e_{i+1}-e_i+c_jt=0.
                                                        \tag{D.IC9}
\]
Eliminating them leaves one free $e_j$ for each orbit,
and the remaining relations are $l_jc_jt=0$.
Hence this quotient is
$\mathbb Z^\rho\oplus\mathbb Z/h$.
Its remaining boundary sends $e_j$ to $c_j\delta$
and sends $t$ to zero. Its kernel is therefore
$\mathbb Z^{\rho-1}\oplus\mathbb Z/h$.
This proves the full nonsplit extension in $H_2$.
For example, two interchanged components of degree
$2$ contribute a relation $4t=0$, not two unrelated
order-two summands. The fibre product torus
$\delta\wedge\gamma$ is the cone generator $t$ itself.

The degree-three cone has the direct sum of the
coinvariants of the component three-cells, of rank
$\rho$, and the invariant component two-cycles with
Chern evaluation zero, of rank $\rho-1$.
Both groups are free. Degree four is the invariant
component three-cell group of rank $\rho$.
The same calculation with $c=0$ gives (D.IC5).

There is also an intrinsic check of the extension.
Let $Y$ be the mapping torus of $K$ and $Q$.
The cap is the principal $S^1_\delta$ bundle over $Y$.
The Euler class evaluates as $c_j$ on the orbit
representatives in $H_2(Y)$, and as $l_jc_j$ on
invariant orbit sums in $H_3(Y)$ paired with the
positive $\gamma$. The integral Gysin sequence
therefore injects $\mathbb Z/h$ into $H_2(U)$ by
$\gamma\mapsto\delta\wedge\gamma$.
The quotient is the free kernel of the first Chern
functional on $H_2(Y)$. This verifies both the
extension and the actual generator without a
choice of an abstract splitting.

\emph{The regular area and the full boundary image.}
In the single circle bundle $P_c\to K$ the map
$H_2(P_c)\to H_2(K)$ is injective, with image
$\ker(c^t)$. The regular elliptic area projects to
$\mu$, the multiplicity cycle of the central fibre,
and hence has its unique lift $(\mu,0)$.
The elliptic loops crossed with the scalar circle
vanish because $H_1(K)=0$.
After passing to $Y$, this area projects to the
orbit vector $(l_j\mu_j)$. Its content is $q_Q$.
The equality with the order of the component of
$Q$ follows directly from the marked component
permutations and multiplicities for the seven listed
types; the retained exact certificate checks every
component class. The actual section-cocycle model
has the same order in $\operatorname{coker}(B-I)$,
so the logarithm $z$ has exactly this denominator.

To identify the rational mixed area coefficients,
pass to the $q_Q$-fold long-circle cover.
The section $q_QQ$ is narrow; its marked local
logarithm straightens the translation through the
original elliptic neighbourhood, as in D.2.16 and
D.2.31. Its corrected long direction is
$q_Q(\gamma-z)$. An elliptic loop crossed with this
direction is zero in the base mapping torus because
the elliptic loop bounds in the Kodaira neighbourhood.
Thus the two mixed free coefficients are $z_2$ and
$-z_1$ times the elliptic-area coefficient.
The integers $q_Q,q_Qz_1,q_Qz_2$ have gcd one.
Consequently these actual fibre tori generate the
primitive free orbit vector
$(l_j\mu_j)/q_Q$.
The torus $\delta\wedge\gamma$ already generates
all cap torsion. This proves the claimed fibre
image and its free cokernel.

Finally $H_2(U,\partial U)=H^4(U)=\mathbb Z^\rho$,
since $H_3(U)$ is free. The boundary first Betti
number is three, while that of $U$ is one for
$c\ne0$ and two for $c=0$. The boundary maps onto
$H_1(U)$: it contains the scalar and long generators.
The relative sequence therefore makes the full
boundary degree-two image have rational rank one
or two, respectively. Those are already the ranks
of the primitive fibre images just proved. Their
cokernels are free, so no further finite enlargement
is possible. The images are equal integrally.
The Wang projection of the full kernel is therefore
all of the fixed lattice.

For completeness, the scalar fixed-normal torus
bounds the scalar-circle bundle over the zero-section
disc. The other fixed-normal torus has zero projection
to $H_2(Y)$: after the same narrow cover and marked
straightening, it is the boundary of the zero-section
disc crossed with the long circle. Since $H_2(Y)$
is free, it was zero before the cover as well.
Its image in $H_2(U)$ is consequently torsion.
If $h=1$ both normal images vanish, and the full
fibre map is the primitive integral area functional
in (D.IC7). This proves the complete marked assertion
in that subsector.
\end{proof}
